
\begin{proposition}[The complete all-ruled cycle criterion with its full elliptic translation]\label{cdp:st:C.2.UN.55}
Specify an oriented cycle of $N\geq1$ marked ruled elliptic surfaces
$R_i=\PP(\OO_{E_i}\oplus L_i)$, with
$L_i\in\operatorname{Pic}^0(E_i)$, elliptic identifications of
successive marked sections, and duality of their actual normal
lines. For $N=1$ identify the two distinct marked sections of
the same normalization. The following are equivalent:
the integral complement-lattice product $M$ of \ref{cdp:st:C.2.UN.51} is the
identity; the specified support has an actual smooth-total
two-torus smoothing; and it occurs as a full reduced fibre of
a smooth compact complex threefold over $\PP^1$ with two-torus
general fibre.

The compact realization retains every marked surface and every
elliptic gluing map, including an arbitrary total elliptic
translation around the cycle. It chooses the normal-product
constants. Its only singular fibres are the specified cycle and
another reduced cycle of $N$ ruled elliptic components. Every
reduced fibre support has Euler number zero.
\end{proposition}

\begin{proof}
The necessary condition and the actual local equivalence are
\ref{cdp:st:C.2.UN.51}. We prove the compact assertion without removing the
horizontal affine translation.

\par\medskip\noindent\textbf{An actual common lattice and its retained translation.}
Use the actual logarithmic linear identifications of \ref{cdp:st:C.2.UN.51} to put
all complements in $G=\C^2/\Lambda$, where $\Lambda$ has real
rank three. All positive outgoing normal meridians are the same
primitive $n\in\Lambda$. Put $p=-n$. The normal-line quotient of
each ruled surface is an elliptic curve. Consequently
$\C p\cap W=\R p$, where $W=\operatorname{span}_{\R}\Lambda$.
The complex line $W\cap iW$ and $p$ are a complex basis of
$\C^2$. The complex linear map sending this line to the first
coordinate line and $p$ to $(0,1)$ puts $W$ exactly in
$\C\times\R$.

Extend the primitive $p$ to an integral basis of $\Lambda$.
After ordering the other two vectors, its exact form is
\[
 \begin{gathered}
 w_1=(a_1,c_1),\qquad w_2=(a_2,c_2),\qquad p=(0,1),\\
 c_1,c_2\in\R,\qquad
 C=\operatorname{Im}(\overline{a_1}a_2)>0.
 \end{gathered}
 \tag{C.TL.1}\label{cdp:eq:C.TL.1}
\]
This changes coordinates, not the lattice or its integral
subgroups. In common coordinates every actual edge map is a
translation with its normal logarithmic term. Choose logarithmic
lifts of its constants $b_i$ and put $b=\sum_i b_i$.
Its first coordinate, modulo
$\Z a_1+\Z a_2$, is exactly the total elliptic gluing translation.
A normal-product constant changes $b_i$ only in the line
$\C p$. Adjust one such constant so that
\[
                         b=(\sigma,i).
 \tag{C.TL.2}\label{cdp:eq:C.TL.2}
\]
The prescribed elliptic maps, hence the horizontal class of
$\sigma$, are unchanged. No restriction on that class has been
imposed.

\par\medskip\noindent\textbf{A holomorphic period family on the entire punctured sphere.}
For $z\in\C$ and a fixed vector $w=(w^{(1)},w^{(2)})\in\C^2$
define
\[
 V_z(w)=
 \bigl(w^{(1)}+z\overline{w^{(2)}},
       w^{(2)}-z\overline{w^{(1)}}\bigr).
 \tag{C.TL.3}\label{cdp:eq:C.TL.3}
\]
For each fixed $w$ this is a holomorphic vector-valued function
of $z$. For each fixed $z$ it is a real linear isomorphism,
with real determinant $(1+|z|^2)^2$. On a branch of $\log z$,
put
\[
 \begin{gathered}
 P_1(z)=V_z(w_1),\quad P_2(z)=V_z(w_2),\quad
 P_3(z)=V_z(p),\\
 P_4(z)=V_z(b)+\frac{N}{2\pi i}\log z\,V_z(p).
 \end{gathered}
 \tag{C.TL.4}\label{cdp:eq:C.TL.4}
\]
Changing the logarithm by $2\pi i$ replaces $P_4$ by
$P_4+NP_3$. Thus the subgroup generated by these four vectors
is single-valued. We check its real rank everywhere, including
where the rank-three group alone has no elliptic normal quotient.

Write $r=|z|$ and $\tau=N\log z/(2\pi i)$. The actual inverse
of \eqref{cdp:eq:C.TL.3} gives
\[
 V_z^{-1}\bigl(iV_z(p)\bigr)
       =\frac{i(2z,1-r^2)}{1+r^2}.
 \tag{C.TL.5}\label{cdp:eq:C.TL.5}
\]
The first three vectors in \eqref{cdp:eq:C.TL.1} span $\C\times\R$, so
their real determinant with any fourth vector $v$ is
$C\operatorname{Im}v^{(2)}$. Equation \eqref{cdp:eq:C.TL.2} therefore gives
$\det_{\R}(w_1,w_2,p,b)=C$. The real multiple
$\operatorname{Re}\tau\,V_z(p)$ has no effect on a determinant.
Using \eqref{cdp:eq:C.TL.5} for the remaining multiple yields the exact value
\[
 \begin{aligned}
 D(z)&=\det_{\R}(P_1(z),P_2(z),P_3(z),P_4(z))\\
 &=C(1+r^2)
       \left(1+r^2+\frac{N}{2\pi}(r^2-1)\log r\right)>0
                         \qquad (0<r<\infty).
 \end{aligned}
 \tag{C.TL.6}\label{cdp:eq:C.TL.6}
\]
Indeed $(r^2-1)\log r\geq0$, while $1+r^2>0$.
In particular the retained horizontal $\sigma$ and the real
characters $c_1,c_2$ require no inequality.

The four holomorphic period sections on a simply connected
base chart act on $\C^2$ by translations. Their real determinant
is nonzero; on each compact base subset their real matrices and
inverses are uniformly bounded. Thus their integral translation
action is properly discontinuous, and a continuously varying
closed real parallelepiped gives properness of the quotient.
The integral change $P_4\mapsto P_4+NP_3$ identifies these
quotients on overlaps. This constructs a smooth proper family
$\mathcal T\to\C^*$ of connected two-tori, with these actual
periods.

\par\medskip\noindent\textbf{The prescribed filling at zero.}
For $z$ sufficiently close to zero, the line
$\C V_z(p)$ is transverse to the real rank-three lattice in
the same sense as at zero. Indeed its transverse determinant
is $C(1+r^2)(1-r^2)$, which is nonzero for $r<1$.
Use the separate ruled moving frame of (C.MX.10):
the normal vector $V_z(n)$ and a transverse integral lattice
lift define an actual elliptic family and a holomorphic
degree-zero line on it. Its projective completion gives the
marked ruled family through each prescribed $R_i$.
The original integral identifications and all finite changes
of logarithm are retained.

Choose the moving edge translation to be $V_z(b_i)$ in the
common coordinates. It induces the actual translation on the
elliptic quotient and the reciprocal normal-line map; these
extend over the marked sections by their line-bundle formulas.
The normal logarithmic coefficient of each edge is
$-V_z(n)/(2\pi i)=V_z(p)/(2\pi i)$. Consequently the entire
cycle has precisely the fourth period in \eqref{cdp:eq:C.TL.4}, not only
its leading asymptotic term. The central value of each edge
is its prescribed elliptic gluing map.

Apply the proper normal-product plumbing of \ref{cdp:st:C.2.UN.51} to these
moving line charts, with actual equation $xy=z$ in each neck.
The moving component cores are smooth over the disc; the
normal-product equation is smooth at $x=y=z=0$.
Uniform flat norms and finitely many bounded cores and necks
give the same Hausdorff proper local filling as in that proof.
Its full central divisor is reduced, and its support is the
specified marked cycle.

On a punctured base chart, send each core's logarithmic
coordinate to the common vector space, adding the successive
edge translations. Across each neck these maps agree modulo
the periods in \eqref{cdp:eq:C.TL.4}. Across a change of branch of $\log z$
their changes are integral multiples of $P_3$. Going around
the cycle adds exactly $P_4$. They therefore give a holomorphic
map to $\mathcal T$. Its fibre maps are the closed-coframe
isomorphisms of \ref{cdp:st:C.2.UN.51}: all four fundamental-group generators
are exactly the displayed periods. Their local inverses
depend holomorphically on the base, so this is an actual
biholomorphism of punctured families.

\par\medskip\noindent\textbf{The second proper filling at infinity.}
Put $s=z^{-1}$ and change the two-dimensional vector-bundle
frame by $\eta_\infty=\eta/z$. In this frame \eqref{cdp:eq:C.TL.3} becomes
\[
 V^\infty_s(w)
       =\bigl(\overline{w^{(2)}},-\overline{w^{(1)}}\bigr)
                          +s w.
 \tag{C.TL.7}\label{cdp:eq:C.TL.7}
\]
This is an actual holomorphic frame change on $\C^*$, and
$V^\infty_0$ is an invertible antilinear map. The periods
in the new frame are exactly
\[
 \begin{gathered}
 V^\infty_s(w_1),\quad V^\infty_s(w_2),\quad V^\infty_s(p),\\
 V^\infty_s(b)
          -\frac{N}{2\pi i}\log s\,V^\infty_s(p).
 \end{gathered}
 \tag{C.TL.8}\label{cdp:eq:C.TL.8}
\]
The ambiguity of the logarithm again adds an integral period.

Near $s=0$ the rank-three lattice and the complex line through
$V^\infty_s(p)$ give a genuine elliptic quotient; the
transversality at zero follows by applying the invertible
antilinear map in \eqref{cdp:eq:C.TL.7} to the original quotient.
Construct $N$ marked ruled components from this actual
nonzero line bundle. Give each positive outgoing meridian
the vector $n_\infty(s)=V^\infty_s(p)$, and each incoming
meridian its negative. Choose the sum of their full affine
edge constants to be $V^\infty_s(b)$, for example by putting
this constant on one edge and zero on the others.
Every such constant induces an actual elliptic translation
and its compatible normal-line map.

The normal-period sum of this opposite-oriented cycle is
$-N V^\infty_s(p)$. Its normal-product plumbing with equation
$xy=s$ therefore has precisely the periods \eqref{cdp:eq:C.TL.8}.
The same finite-core construction is a smooth proper local
filling with reduced central divisor. Its exact logarithmic
maps identify the punctured family with the frame change of
$\mathcal T$ just displayed. This explicitly constructs the
second filling; no contraction of a global auxiliary line
bundle is involved.

\par\medskip\noindent\textbf{Compact gluing and the Euler comparison.}
Choose small disjoint discs about $0$ and $\infty$ and a middle
annulus covering the rest of $\PP^1$. Glue their two proper
fillings to $\mathcal T$ by the actual punctured-family
biholomorphisms above. Over every point of the base all pieces
are identified with one Hausdorff proper local model. Hence
the glued space is Hausdorff and the map is proper locally
on the base. A finite closed refinement of the three base
charts makes its inverse image a finite union of compact
sets, proving properness over the compact $\PP^1$.
The total space is smooth by its displayed local models.
All other fibres are the tori of \eqref{cdp:eq:C.TL.4}.

Each ruled elliptic component has Euler number
$e(\PP^1)e(E)=0$, and each double curve has Euler number zero.
There are no triple points, so both singular supports have
Euler number zero. This proves every asserted compact and
marked-data comparison, including the one-component
self-identification.

The symbolic determinant, frame-change, inverse-map and
monodromy identities are certified in
\path{code/A_twistor_log_cycles.py} and its saved output.
\end{proof}


